\section{Proofs of KoalaFalcon Requirements}\label{appendix-detailed-loss-term-derivations}

\subsection{Algebraic Feasibility Proofs}\label{appendix-alg-feas}

\paragraph{Proof 1: Ring Algebraic Requirements}

Let $q = 2^{31} - 2^{24} + 1$, a known prime. The field $\mathbb{F}_q$ has multiplicative group $\mathbb{F}_q^\times$ of cyclic order $q-1$. There exists a primitive $2n$-th root of unity in $\mathbb{F}_q$ if and only if $2n | q - 1$. 

$$
q - 1 = 2^{31} - 2^{24}  = 2^{24}(2^7 - 1) 
$$

For $n=512=2^9$, $2^{24}(2^7 - 1) \equiv 0  \mod 2^{10}$. For $n=1024=2^{10}$, $2^{24}(2^7 - 1) \equiv 0  \mod 2^{11}$. Therefore there exists a 
primitive $2n$-th root of unity in $\mathbb{F}_q$ for both choices of $n$. The existence of a primitive $2n$-th root of unity implies that 
$X^n + 1$ has $n$ roots, splitting into $n$ linear factors. $\blacksquare$

\paragraph{Proof 2: Public, Private Key Remain Bases to the Same Lattice}

Let $\Lambda = \Lambda(B_1)$ be a $2n$-dimensional lattice with basis $B_1$. $B_2$ is also a basis of $\Lambda$ if and only if there exists a unimodular matrix $U$ with determinant $\pm 1$ such that $UB_2 = B_1$. This definition is adapted from Theorem~2.9 of \cite{Menezes2026GentleIntro}.

Let $q = 2^{31} - 2^{24} + 1$, $n=2^\kappa$ is a power-of-two, $R = \mathbb{Z}[X]/(X^n + 1)$, with $\phi = X^n + 1$ fully splits into $n$ linear factors over $\mathbb{Z}_q$, and $f, g \in R_q^\times$. The relation

$$
\begin{aligned}
h &\equiv g f^{-1} \pmod{\phi, q}, \\
fG - gF &\equiv q \pmod{\phi}
\end{aligned}
$$

The private basis $B_1$ and public basis $B_2$ of a $2n$ dimensional NTRU lattice are defined below. 

$$
B_1 =
\begin{bmatrix}
g & -f \\
G & -F
\end{bmatrix},
\qquad
B_2 =
\begin{bmatrix}
-h & I_n \\
qI_n & O_n
\end{bmatrix}.
$$

We wish to show there exists a unimodular matrix $U$ with integral coefficients such that $UB_2 = B_1$. One such matrix is 

$$
U =
\begin{bmatrix}
-f & \frac{g-hf}{q} \\
-F & \frac{G-hF}{q}
\end{bmatrix}.
$$

% I believe this is a standard example of a unimodular matrix for transforming the private basis to the public basis, should be cited at some point

Evaluating $UB_2$ gives

$$
\begin{aligned}
U B_2
&=
\begin{bmatrix} -f & \frac{g-hf}{q} \\ -F & \frac{G-hF}{q} \end{bmatrix}
\begin{bmatrix} -h & I_n \\ qI_n & O_n \end{bmatrix}
\\[4pt]
&=
\begin{bmatrix}
-f(-h) + \frac{g-hf}{q}\, qI_n & (-f)I_n + \frac{g-hf}{q}\, O_n \\[4pt]
-F(-h) + \frac{G-hF}{q}\, qI_n & -F\, I_n + \frac{G-hF}{q}\, O_n
\end{bmatrix}
\\[4pt]
&=
\begin{bmatrix}
fh + g - hf & -f \\
Fh + (G-hF) & -F
\end{bmatrix}
\\[4pt]
&=
\begin{bmatrix} g & -f \\ G & -F \end{bmatrix}
= B_1 .
\end{aligned}
$$

The determinant of $U$ is 

$$
\det(U) = (-f \cdot \frac{G - hF}{q}) - (-F \cdot \frac{g - hf}{q}) 
$$
$$
= \frac{-fG + fhF}{q} - \frac{-Fg + Fhf}{q} 
= \frac{-fG + fhF + Fg - Fhf}{q} = \frac{-fG + Fg}{q}
$$

And since $fG - gF \equiv q \pmod{X^n+1}$,

$$
\det(U) = \frac{-fG + Fg}{q} = - \frac{q}{q} = -1
$$

Therefore there exists a unimodular matrix $U$ with integral coefficients such that $UB_2 = B_1$. Thus $B_1, B_2$ are bases of the same lattice $\Lambda$. $\blacksquare$

% TODO: maybe justify why coefficients are integral, this is probably fine as is but technically the proof is incomplete until then.

\paragraph{Proof 3: Standard Deviation is Greater than Smoothing Parameter}\label{proof-3}

The loss-term corollaries require
\(s \ge \eta_\epsilon(\Lambda_{h,q})\). By the GPV smoothing lemma
(Lemma~3.1 of \cite{GPV2008}), for the trapdoor basis \(B_{f,g}\) produced by
\(\mathrm{TpdGen}\),

\[
\eta_\epsilon(\Lambda_{h,q}) \le \eta_\epsilon(\mathbb{Z}^{2n}) \cdot ||B_{f,g}||_{GS}.
\]

By Lemma~4 of \cite{Fouque2024CoreFalcon},
\(\eta_\epsilon(\mathbb{Z}^{2n}) \le \eta_\epsilon^{\mathrm{upper}}(\mathbb{Z}^{2n})\).
With \(||B_{f,g}||_{GS} \le \alpha\sqrt{q}\) and the CoreFalcon+ choice
of \(s\),

\[
\eta_\epsilon(\Lambda_{h,q})
\le
\eta_\epsilon^{\mathrm{upper}}(\mathbb{Z}^{2n}) \cdot \alpha\sqrt{q}
= s.
\]

$\blacksquare$

For intuitive context, the trapdoor-quality bound
\(||B_{f,g}||_{GS} \le \alpha\sqrt{q}\) forces the relevant smoothing
threshold to scale as \(\sqrt{q}\), and the definition of \(s\) matches
it. Consider the integer lattice \(\mathbb{Z}^{2n}\) scaled by
\(||B_{f,g}||_{GS}\). Since \(\Lambda_{h,q}\) may be skewed, the scaled
integer lattice \(||B_{f,g}||_{GS}\cdot \mathbb{Z}^{2n}\) represents the
case where every orthogonal direction of a \(2n\)-dimensional lattice is
of length equal to the worst-case Gram-Schmidt norm of
\(||B_{f,g}||_{GS}\). Smoothing this scaled lattice requires width at
least \(||B_{f,g}||_{GS} \cdot \eta_\epsilon(\mathbb{Z}^{2n})\), which
is exactly the CoreFalcon+ choice of \(s\) up to the Lemma 4 upper
bound. Raising \(q\) raises both sides by the same \(\sqrt{q}\) factor,
so the condition is preserved.

